% Sección 2
\section{Qué pasó}

\subsection{El sismo y la emergencia}

El sismo ocurrió el 10 de agosto de 2026 a las 7:34 de la mañana. Tuvo magnitud 7,4 y epicentro en San José del Palmar, en el Chocó, según el boletín del Servicio Geológico Colombiano que citan los considerandos del Decreto 1261. El mismo decreto da dos profundidades distintas, 103,4 km en un considerando y unos 94 km en otro. El 19 de agosto, el Gobierno declaró el estado de emergencia por 30 días calendario en quince departamentos, entre ellos Caldas.

El balance nacional creció con los días. El decreto cita el reporte 36 de la UNGRD, con corte al 19 de agosto. Ese reporte contaba 314 personas fallecidas, 4.437 heridas y 262 desaparecidas, y 262.154 damnificados en 147.464 familias. Seis días después, la Organización Panamericana de la Salud (OPS) citaba otro corte de la UNGRD, con 331 fallecidos y 361.707 afectados. Las cifras tempranas eran más bajas. El 10 de agosto Asocapitales contaba 87 muertos en las capitales, y el 12 de agosto la UNGRD reportaba 241, según la Federación Internacional de la Cruz Roja (IFRC). En octubre, la Defensoría hablaba de 335 muertos en 500 municipios de 16 departamentos.

La cobertura territorial también cambió de un corte a otro. El decreto habla de 471 municipios con afectación directa en quince departamentos, al 19 de agosto. La OPS, con datos de la UNGRD al 25 de agosto, cuenta 494 municipios en 16 departamentos.

Las vías también resultaron afectadas. Al 19 de agosto, la UNGRD contaba 449 vías afectadas, cinco aeropuertos y varios puentes. El 12 de agosto, según la IFRC, la cifra era de 125 vías dañadas y dos aeropuertos.

La energía ayuda a entender lo que pasó con los perecederos. Según el decreto, la demanda de energía no atendida en la región de Caldas, Quindío y Risaralda fue al comienzo del 82~\%, y el 13 de agosto había bajado al 5,9~\%. El aeropuerto La Nubia, en Manizales, también resultó afectado.

El decreto mide además el peso de la zona afectada en la comida del país. Según el SIPSA del DANE, los departamentos afectados aportan el 11,5~\% del abastecimiento nacional de alimentos, y el 30,8~\% en el caso de las carnes. El decreto advierte presión sobre los precios por las interrupciones del transporte. También señala que la continuidad del programa de alimentación escolar quedó sin garantía.

\subsection{Manizales y Caldas}

Las cifras de damnificados de Manizales cambiaron con la fecha y con el método de conteo. El 11 de agosto, el Presidente habló en la ciudad de unos 2.000 damnificados y 1.150 familias sin vivienda. El 13 de agosto, la Alcaldía decía haber censado a más de 2.000 personas. Los reportes de la Alcaldía y de su Unidad de Gestión del Riesgo (UGR) que recoge Asocapitales hablan de 4.000 damnificados, con cortes del 13 y el 15 de agosto. El censo que cerró el 24 de agosto registró 8.240 personas damnificadas y 3.080 predios evacuados, según el alcalde Jorge Eduardo Rojas. La misma nota informa que siete edificios con daño estructural grave serán demolidos.

Las muertes en la ciudad también tienen dos cifras. El primer balance de la Alcaldía, del 10 de agosto, contaba cuatro. Los reportes con corte al 13 y al 15 de agosto subieron la cifra a seis. Para Caldas hay dos versiones, cinco muertos según la UNGRD citada por la OPS al 25 de agosto y seis según el vicepresidente José Manuel Restrepo el 18 de agosto.

Las viviendas dañadas de Manizales aparecen con fechas que no cuadran entre los informes del equipo. Según la investigación de evidencia, Asocapitales registra 1.449 viviendas colapsadas y 1.970 averiadas con corte al 13 de agosto, y 2.707 colapsadas y 3.993 averiadas con corte al 15 de agosto. La investigación de referentes cita otro informe de Asocapitales que asigna las cifras de 2.707 y 3.993 al corte del 13 de agosto. La fecha de esas dos cifras queda \porverificar.

En Caldas, la OPS cita a la UNGRD con 27 municipios afectados, 9.511 personas afectadas, 329 heridos y tres albergues temporales con 230 personas, al 25 de agosto. Las viviendas dañadas del departamento tienen tres versiones que no coinciden. El 12 de agosto, según la IFRC, la UNGRD contaba 4.064 viviendas averiadas y 1.442 destruidas. El 18 de agosto, el vicepresidente habló de 9.400 viviendas con afectación parcial y 2.200 con afectación total. La Gobernación, con datos de solo ocho municipios, reportó el 20 de agosto 2.658 con pérdida parcial y 582 con pérdida total.

Los coliseos de la ciudad se volvieron albergue y centro de acopio. El 10 de agosto, el Coliseo Mayor y el Coliseo Menor recibieron a personas evacuadas. El albergue junto al Coliseo Menor pasó de 114 a 245 personas el 12 de agosto, y el Coliseo Mayor tenía unas 300 personas el 13 de agosto. El 24 de agosto el Coliseo Mayor seguía siendo el albergue principal. Esa cifra no encaja con la de la OPS, que el 25 de agosto contaba 230 personas en tres albergues temporales de todo Caldas. Una explicación posible, que es inferencia nuestra, es que las dos fuentes cuenten tipos distintos de albergue. El dato queda \porverificar.

La respuesta del primer día movilizó más de 324 unidades de rescate y casi 2.000 policías. El alcalde pidió no sacar vehículos a la calle para dejar paso a ambulancias y bomberos. El censo de damnificados se hizo en cuatro puntos de la ciudad, que fueron Milán, Chipre, la avenida Santander y el centro. Lo coordinaron Bomberos, Defensa Civil y Cruz Roja con la Alcaldía. Estar en ese censo era requisito para recibir el subsidio de arriendo, de 300.000 pesos mensuales por tres meses prorrogables. El comercio del centro estuvo cerrado nueve días, y la Alcaldía autorizó su reapertura gradual el 18 de agosto. Con corte al 13 de agosto, la ciudad pedía colchonetas, mantas, carpas y plantas eléctricas para el albergue.

\subsection{El corredor de los bancos de alimentos}

La ayuda alimentaria tuvo en Manizales uno de sus nodos. El 10 de agosto, la red de bancos de alimentos del país (ABACO) activó un corredor humanitario, con Armenia, Cali e Ibagué como puntos intermedios. Los bancos de Pereira y Buenaventura tuvieron que buscar alternativas por daños en sus sedes. La red anunció que daría prioridad a los albergues oficiales, los hospitales y las familias autoalbergadas.

El papel de Manizales en ese corredor tiene dos lecturas. Las dos investigaciones del equipo resumen la misma nota de Noticias RCN de forma distinta. Una dice que Manizales fue el punto principal de acopio, y la otra, que fue uno de los principales. La diferencia queda \porverificar{} abriendo la nota.

El volumen del corredor se conoce solo para toda la red. El 18 de agosto, ABACO había superado las 2.000 toneladas y focalizado a más de 9.000 familias, con 23 bancos activos. Al cierre de agosto hablaba de más de 9.470 familias en 36 municipios. Esas cifras no dicen cuánto se repartió en Manizales.

La misma red nombró sus dificultades el 18 de agosto, por boca de su gerente de Cooperación Internacional. El reto principal fue la llegada al territorio. También mencionó el transporte disponible, la articulación entre actores, el almacenamiento en los acopios y la falta de vías de acceso claras, sobre todo hacia las zonas rurales. Para entonces había tres bancos de alimentos afectados, en Buenaventura, Cartago y Pereira, y el de Pereira podría tener que demolerse.

Los perecederos quedaron en segundo plano. Según el director ejecutivo de ABACO, los cortes de energía impedían manejar alimentos que necesitan frío, y por eso la red priorizó enlatados y alimentos de larga vida. La Cruz Roja, por su parte, pidió el 14 de agosto no donar perecederos, alimentos abiertos, agua con empaque dañado ni medicamentos abiertos o por vencer. Ese día informó que había despachado más de 280 toneladas a Buenaventura, Cali, Quibdó y Pereira.

Lo que pasó en Bogotá muestra dónde se forma el cuello de botella. El centro de acopio del estadio El Campín reunió 933 toneladas en seis días. Al cierre seguía pidiendo voluntarios, montacargas y camiones para clasificar y despachar lo recibido. De esa carga, 157 toneladas le correspondían a Caldas. La investigación no registra la fecha de esa nota.

En Manizales, el Coliseo Menor funcionó como centro de acopio, coordinado por la gestora social del municipio, Juliana Londoño. El 12 de agosto recibía mercados, artículos de aseo, cobijas y colchonetas. Según ella, repartía entre más de 4.000 desplazados. Llegaron además otras ayudas, sin cifras para la ciudad. El 11 de agosto el Presidente anunció kits de aseo, alimentación, noche y cocina para Manizales. El 15 de agosto llegó una tractomula con ayuda de la Alcaldía de Bogotá. En septiembre, la prensa reportó entregas de kits de mercado de una fundación privada en Manizales, Supía y Riosucio, dentro de un total de 8.000 kits.

El abastecimiento mayorista de la ciudad bajó la semana siguiente al sismo. Según el SIPSA del DANE, la plaza Centro Galerías recibió 2.069 toneladas entre el 13 y el 19 de agosto, un 4,46~\% menos que las 2.165 de la semana anterior. Los tubérculos cayeron un 20,73~\% y las verduras un 10,02~\%. El boletín atribuye la caída nacional a problemas de recolección después del sismo, pero no vincula de forma directa el dato de Manizales con el sismo.

Las quejas documentadas hasta ahora tratan de ayudas económicas y no de comida. El 5 de octubre, la Procuraduría dijo que gobernadores y alcaldes de Risaralda, Quindío y Caldas señalaban desembolsos tardíos y una distribución desigual. El Registro Único de Damnificados cerró el 4 de septiembre y dejó por fuera a familias rurales dispersas y a familias étnicas.

\subsection{Las vías cerradas}

Dos cierres totales afectaron las salidas de la ciudad en los primeros días. El primero estaba en la vía Pereira–Chinchiná, entre los kilómetros 22 y 24. La Policía de Tránsito lo reportó el 10 de agosto, y seguía vigente el 12 de agosto a las 9:24 de la mañana y a las 7:10 de la noche. La prensa dice que el cierre afecta la conexión del sur de Caldas con Pereira. Los documentos del equipo lo cuentan como una de las salidas de Manizales. Si el tramo entre Manizales y Chinchiná estuvo cerrado queda \porverificar.

El segundo cierre estaba en la vía Manizales–Fresno, en el Alto de Letras, entre los kilómetros 36 y 40, ya en Herveo (Tolima). Bloqueaba la salida de Manizales hacia Fresno, Mariquita y Honda, y seguía vigente el 12 de agosto. Según Portafolio, la UNGRD todavía incluía los dos corredores entre los afectados en su corte del 19 de agosto.

El corredor concesionado Armenia–Pereira–Manizales tuvo derrumbes sobre la calzada en varios puntos, como Tarapacá I y II, Club Campestre–Caimo y La Paila–Río La Vieja. La Agencia Nacional de Infraestructura (ANI) lo informó el 10 de agosto. Según El Colombiano, que cita al INVÍAS, los daños típicos fueron deslizamientos traslacionales, caída de material vegetal y daños totales o parciales en las calzadas. El 25 de agosto, la OPS seguía reportando restricciones de acceso terrestre al Eje Cafetero.

Varias preguntas sobre las vías siguen sin fuente. No encontramos reportes de cierres por el sismo en la vía de La Siria ni en las salidas a Medellín y a Bogotá. La vía a Bogotá, en cambio, pasa por Letras y Fresno, que sí figuran cerradas. Tampoco encontramos la fecha de reapertura de ninguno de los dos cierres, ni reportes de embotellamientos en los accesos a la ciudad.

\subsection{Las cifras en una tabla}

La tabla~\ref{tab:cifras} reúne las cifras de esta sección con su fecha de corte, su fuente y su nivel de evidencia. Las fuentes completas, con enlace, están en el anexo.

{\footnotesize
\setlength{\tabcolsep}{4pt}
\renewcommand{\arraystretch}{1.25}
\begin{longtable}{L{0.17\linewidth} L{0.31\linewidth} L{0.13\linewidth} L{0.18\linewidth} L{0.12\linewidth}}
\caption{Cifras del sismo con su fecha de corte. Ninguna ha sido abierta todavía por una persona del equipo.}\label{tab:cifras}\\
\toprule
\textbf{Dato} & \textbf{Valor} & \textbf{Corte} & \textbf{Fuente} & \textbf{Nivel}\\
\midrule
\endfirsthead
\multicolumn{5}{l}{\small\itshape Tabla~\ref{tab:cifras}, continuación}\\
\toprule
\textbf{Dato} & \textbf{Valor} & \textbf{Corte} & \textbf{Fuente} & \textbf{Nivel}\\
\midrule
\endhead
\midrule
\multicolumn{5}{r}{\small\itshape Sigue en la página siguiente}\\
\endfoot
\bottomrule
\endlastfoot
Magnitud, hora y epicentro & 7,4 Mw; 7:34 a.\,m.; San José del Palmar (Chocó) & 10 ago 2026 & Decreto 1261, que cita al SGC & \textbf{Oficial} (alta)\\
Profundidad & 103,4 km en un considerando y unos 94 km en otro & 19 ago 2026 & Decreto 1261 & \textbf{Oficial} (media)\\
Emergencia & 30 días calendario en 15 departamentos, Caldas entre ellos & 19 ago 2026 & Decreto 1261 & \textbf{Oficial} (alta)\\
Balance nacional (reporte 36) & 314 fallecidos, 4.437 heridos y 262 desaparecidos; 262.154 damnificados en 147.464 familias & 19 ago 2026 & Decreto 1261, que cita a la UNGRD & \textbf{Oficial} (alta, preliminar)\\
Balance nacional posterior & 331 fallecidos, 4.439 heridos, 240 desaparecidos y 361.707 afectados & 25 ago 2026 & OPS, informe de situación 6, que cita a la UNGRD & \textbf{Oficial} (alta)\\
Muertos, otros cortes & 87 en capitales (10 ago); 241 (12 ago); 335 en 500 municipios (octubre) & 10 ago a oct 2026 & Asocapitales n.º 6; IFRC; Defensoría vía Minuto60 & \textbf{Prensa} (sin calificar)\\
Municipios afectados & 471 en 15 departamentos (19 ago); 494 en 16 departamentos (25 ago) & 19 y 25 ago 2026 & Decreto 1261; OPS & \textbf{Oficial} (alta)\\
Vías afectadas en el país & 449 vías, 5 aeropuertos y varios puentes (19 ago); 125 vías y 2 aeropuertos (12 ago) & 12 y 19 ago 2026 & Decreto 1261; IFRC & \textbf{Oficial} (alta); \textbf{Prensa} (sin calificar)\\
Energía en Caldas, Quindío y Risaralda & Demanda no atendida del 82~\% al comienzo y del 5,9~\% el 13 de agosto & 13 ago 2026 & Decreto 1261 & \textbf{Oficial} (alta)\\
Peso en el abastecimiento & 11,5~\% de los alimentos del país; 30,8~\% de las carnes & 19 ago 2026 & Decreto 1261, que cita al SIPSA-DANE & \textbf{Oficial} (alta)\\
Muertos en Manizales & 4 en el primer balance; 6 en los cortes del 13 y el 15 de agosto & 10, 13 y 15 ago 2026 & Noticias RCN; Asocapitales n.º 22 y n.º 27 & \textbf{Prensa} (media)\\
Muertos en Caldas & 5 (UNGRD vía OPS); 6 (vicepresidente) & 25 y 18 ago 2026 & OPS; Blu Radio & \textbf{Oficial} y \textbf{Prensa} (media)\\
Damnificados en Manizales & Unos 2.000 (11 ago); más de 2.000 censados (13 ago); 4.000 (13 y 15 ago); 8.240 personas y 3.080 predios evacuados (24 ago) & 11 a 24 ago 2026 & NTN24; Blu Radio; Asocapitales; Noticias RCN & \textbf{Prensa} (media)\\
Viviendas en Manizales & 1.449 colapsadas y 1.970 averiadas (13 ago); 2.707 y 3.993 (15 ago, otro informe las fecha el 13 ago: \porverificar) & 13 y 15 ago 2026 & Asocapitales n.º 22, n.º 24 y n.º 27 & \textbf{Prensa} (media)\\
Caldas según la UNGRD & 27 municipios; 9.511 personas afectadas; 329 heridos; 3 albergues temporales con 230 personas & 25 ago 2026 & OPS, que cita a la UNGRD & \textbf{Oficial} (alta)\\
Viviendas en Caldas & 4.064 averiadas y 1.442 destruidas (12 ago); 9.400 parciales y 2.200 totales (18 ago); 2.658 y 582 en 8 municipios (20 ago) & 12 a 20 ago 2026 & IFRC; Blu Radio (vicepresidente); Portafolio (Gobernación) & \textbf{Prensa} (media; IFRC sin calificar)\\
Albergues en Manizales & De 114 a 245 personas junto al Coliseo Menor (12 ago); unas 300 en el Coliseo Mayor (13 ago); Coliseo Mayor como albergue principal (24 ago) & 10 a 24 ago 2026 & Noticias RCN; Blu Radio & \textbf{Prensa} (media)\\
Corredor de ABACO & Más de 2.000 t y más de 9.000 familias (18 ago); más de 9.470 familias en 36 municipios (cierre de agosto) & 18 ago y 21 sep 2026 & Noticias RCN; Portafolio & \textbf{Prensa} (media)\\
Acopio de El Campín (Bogotá) & 933 t en seis días; 157 t para Caldas & Sin fecha en la investigación & Noticias RCN & \textbf{Prensa} (sin calificar)\\
Cierres de salida & Pereira–Chinchiná, km 22 a 24, y Manizales–Fresno, km 36 a 40; ambos totales y vigentes el 12 ago & 10 y 12 ago 2026 & Publimetro; Semana; Colombia.com; Minuto30 & \textbf{Prensa} (media)\\
Plaza Centro Galerías & 2.069 t entre el 13 y el 19 de agosto, 4,46~\% menos que la semana anterior & 15 a 21 ago 2026 & DANE, SIPSA, boletín semanal & \textbf{Oficial} (media)\\
\end{longtable}
}
